\section{Discussion and lessons learned}
\label{sec:discussion}
A document operation is not complete from a retrieval user's perspective
until search reaches the right version and path. The early rename test
found stale entries in the index; the later retest found stale labels only
in a cached response. A search probe checks what the agent can find, while
an index inspection distinguishes these two failure modes. Neither
replaces checks of the files, tree registry and library metadata
\citep{park-metadata-quality-2009,paivarinta-ecm-2005}.

Once search reaches the intended collection, the next question is which
text reaches the reader. The description-bearing cards performed better
on the scripted routing task than the nominally name-only cards, but A0
still missed target strings in correctly selected collections. A0h closed
part of that gap by adding a search lane and changing the pack order. The
retrieval trace counted more text for A0h than for A0, and A0h did not
significantly outperform the similarly costly unrouted PD configuration.
The result favours testing
routing and evidence construction separately rather than crediting the
whole gain to the hierarchy.

The verifier illustrates a related distinction. In fixed pools, sorting by
score pushed some gold documents out of the top five. A separate archive
had no recorded scores below the admission cutoff. It therefore tells us
little about the value of filtering, even though the replay documents an
ordering effect. An operator needs to be able to inspect both decisions.

The native answer comparison allowed HMRS to revisit searches, search
outside its five configured collections and read source text, unlike the
capped one-shot baselines. Its higher score in the
unified pass depended on the judging protocol. This comparison does not
isolate backend performance and did not measure the ability of the full
catalogue-first workflow to support individual answer claims.

Finally, an answer should say which collections were searched. In E7,
search returned success after skipping eight unreadable collections.
The E2 agent shows the same need from the other side: it was configured
for five collections but searched every collection on the host in nine of
ten questions, and only its stored calls reveal this.
Without that information, a user cannot tell whether the system looked
and found nothing, or did not reach part of the corpus.
